\documentclass[11pt]{article}
\usepackage{amsmath,amssymb,booktabs,geometry,hyperref}
\geometry{margin=1in}

\title{Dependency-Preserving Graph Integration and Hilbert Lower Certificates}
\author{intervalNets implementation note}
\date{}

\begin{document}
\maketitle

\section{Scope and semantics}

This note specifies the graph/Hilbert certification path implemented in
\texttt{graph\_hilbert.py}.  It supplements, rather than replaces, the
coefficientwise pointwise-residual functional in
\texttt{certified\_polynomial\_zonotope\_integration.tex}.  Approximation
noise remains pointwise: its value may vary with the domain point.  The new
method does not integrate such a symbol as one global constant.

All formulas below use normalized measure
\[
 d\mu(x)=|\Omega|^{-1}\,dx.
\]
Raw squared integrals are recovered by multiplication with \(|\Omega|\).
For the affine box parameterization \(x=c+G\alpha\), the implemented scalable
backend currently requires an axis-aligned \(G\) for the Neumann witness, while
the exact reference moment backend works on the reference box directly.

\section{Immutable arithmetic graph}

The exact factored Jacobian is translated to immutable nodes for constants,
noise slices, addition, linear maps, Hadamard products, componentwise squares,
and diagonal scaling.  Every node records its domain and approximation-noise
support.  Nodes are hash-consed, so repeated occurrences of one approximation
symbol refer to the same graph input.  For every fixed noise vector,
\[
 \operatorname{Eval}(\mathcal G;\alpha,\eta)
 =\operatorname{Eval}(J_{\rm factored};\alpha,\eta).
\]

The \texttt{SparseReferenceMomentBackend} expands a small domain-only graph to
canonical exponent maps and computes
\[
 \mathbb E_\mu\!\left[\sum_k P_k(\alpha)^2\right]
 =\sum_{k,\beta,\gamma}
 c_{k,\beta}c_{k,\gamma}\,
 \mathbb E[\alpha^{\beta+\gamma}].
\]
It is an exact regression oracle, not the scalable 100-dimensional backend.
Trying to expand a degree-six dense graph in 100 variables would merely hide
the original combinatorial problem.

\section{Certified Hilbert compression}

At every activation boundary, the scalable backend stores
\[
 y(\alpha)=\widehat y(\alpha)+e(\alpha),\qquad
 \widehat y(\alpha)=a_0+A\alpha,qquad
 \|e\|_{L^2_\mu(\ell_2)}\leq\varepsilon.
\]
The discarded function is a single certified Hilbert remainder.  It is not
replaced by independent pointwise symbols.

For an affine preactivation \(q=c+a^\top\alpha\), a certified polynomial
enclosure
\[
 \tanh(q)\in p(q)+[-\delta,\delta]
\]
is projected orthogonally onto
\(\operatorname{span}\{1,\alpha_1,\ldots,\alpha_d\}\).  The implementation
computes the required moments of \(q^k\) without monomial expansion, by
convolving the one-dimensional uniform moments of the independent summands.
The projection remainder is
\[
 \varepsilon_{\rm proj}^2
 =\mathbb E[p(q)^2]-|\Pi_1p(q)|^2_{L^2_\mu},
\]
with an explicit floating-point padding.  Consequently,
\[
 \|\tanh(q)-\Pi_1p(q)\|_{L^2_\mu}
 \leq \varepsilon_{\rm proj}+\delta.
\]
If the incoming preactivation has Hilbert error \(\varepsilon_z\), the
one-Lipschitz property of \(\tanh\) gives the vector bound
\[
 \varepsilon_y\leq \varepsilon_z+
 \left(\sum_i(\varepsilon_{{\rm proj},i}+\delta_i)^2\right)^{1/2}.
\]
Affine layers use a padded spectral-norm upper bound.

For the reverse derivative graph, \(\tanh'\) is treated by its certified
quadratic enclosure.  Products of affine graph nodes are projected without
expansion.  If
\(u=c+a^\top\alpha\) and \(v=d+b^\top\alpha\), their affine projection is
\[
 \Pi_1(uv)=cd+\tfrac13a^\top b+(cb+da)^\top\alpha.
\]
Its exact fourth-moment energy is used to certify the discarded quadratic
part.  The reverse error recurrence uses
\[
 aD-\widehat a\widehat D=(a-\widehat a)D
 +\widehat a(D-\widehat D),
\]
the exact bound \(0\leq\tanh'\leq1\), an affine supremum for
\(\widehat a\), and padded spectral norms at linear maps.

\section{Reverse-triangle certificates}

For either the value field or the stacked value/gradient field, let
\(G=P+R\) with exact affine nominal energy \(A=\|P\|_\mu^2\) and
\(\|R\|_\mu\leq E\).  Then
\[
 \max\{0,\sqrt A-E\}\leq\|G\|_\mu\leq\sqrt A+E.
\]
The result is intersected with the previous absolute-moment/parity interval.
Thus the new upper endpoint cannot increase and the lower endpoint cannot
decrease.

\section{Neumann dual lower bound}

For \(H=W^{1,2}(\Omega)\), any scalar witness \(v\) yields
\[
 \|f\|_H\geq
 \frac{|\langle f,v\rangle_H|}{\|v\|_H}.
\]
The implemented basis is
\[
 \phi_0=1,\qquad
 \phi_i(\alpha)=\frac{\alpha_i^3}{3}-\alpha_i.
\]
It satisfies \(\partial_n\phi_i=0\) on the boundary of an axis-aligned box.
Integration by parts therefore gives
\[
 \langle f,v\rangle_{W^{1,2}_\mu}
 =\int_\Omega f(v-\Delta_xv)\,d\mu.
\]
For \(f=P+R\) and \(\|R\|_{L^2_\mu}\leq E\),
\[
 \|f\|_{W^{1,2}_\mu}\geq
 \frac{\max\{0,|\langle P,v-\Delta v\rangle|-E\|v-\Delta v\|\}}
 {\|v\|_{W^{1,2}_\mu}}.
\]
All Gram matrices are diagonal and analytic for this basis.  A deterministic
one-parameter family of coefficient directions is searched numerically; only
the final candidate is used in the outward-padded certificate.  Optimizer
optimality is never assumed.

\section{Soundness and limitations}

\begin{itemize}
 \item The exact graph backend is used to verify graph evaluation and moment
 equivalence on small cases.
 \item The scalable backend is a certified compression, not exact integration
 of the full 100-dimensional degree-six nominal graph.
 \item Polynomial proposals are not proofs.  Their uniform residuals are
 certified by the repository's outward-rounded subdivision routine.
 \item A larger moment budget or subdivision count is accepted only through
 explicit recomputation; sampled values are diagnostics, never certificates.
 \item The combined result always intersects the pre-existing certificate.
\end{itemize}

The diagnostic decomposition
\[
 \|P\|_\mu\quad\hbox{versus}\quad E
\]
must be reported for both value and gradient fields.  This distinguishes
moment/projection loss from activation-residual and graph-remainder loss.

\end{document}
